\section*{結言}
今日の特徴は、モデルの能力向上だけでなく、能力をどの範囲で提供し、どう統制するかが同時に議論された点にある。Anthropicは利用規約、脆弱性スキャン、研究機関支援を通じて安全性と社会実装の両面を打ち出し、OpenAIとStepFunは高速化や試用経路の拡大を訴えた。

研究投稿は、エージェントの時間管理、共感と事実整合、規則遵守、長文脈の処理効率といった、実運用時の振る舞いを問うものが多かった。今後は性能・価格だけでなく、検証可能な評価条件と利用時の制約が比較の焦点になる。本号の数値や評価はGrokの収集時点の主張を記録したもので、未確認事項は確定情報と区別している。
